\documentclass[11pt,a4paper]{article}

\usepackage[utf8]{inputenc}
\usepackage{amsmath,amsfonts,amssymb}
\usepackage{geometry}
\usepackage{graphicx}
\usepackage{hyperref}
\usepackage{cite}

\geometry{margin=2.5cm}

\title{The REN Formula: Construction, Properties, and Application in Hutchinson Synthesis}
\author{}
\date{}

\begin{document}

\maketitle

\tableofcontents
\newpage

\section{Introduction}

This document presents a family of polar curves named \textbf{REN (Radial Exponential Normalised)}, designed with the following goals:
\begin{itemize}
    \item maximum formal simplicity
    \item automatic normalisation
    \item high geometric expressiveness
\end{itemize}

The formula evolves from a fully symmetric base version to deformed, locally controllable shapes.

\section{The Base REN Formula}

The fundamental formula is:

\begin{equation}
r(\theta) = \exp\left(-\lambda \left| \sin(k\theta) \right|^q \right)
\end{equation}

\subsection{Definition and Parameters}
\begin{itemize}
    \item $k \in \mathbb{N}^+$: number of lobes
    \item $\lambda > 0$: depth of the minima
    \item $q > 0$: shape control (sharpness)
\end{itemize}

\subsection{Properties}
\begin{itemize}
    \item $r_{max} = 1$ (normalisation)
    \item periodicity: $\frac{2\pi}{k}$
    \item full symmetry
\end{itemize}

\subsection{Geometric Interpretation}

The function $\sin(k\theta)$ defines an angular distance from the symmetry points.
The exponential introduces a controlled decay.

\section{Extending REN with Angular Mobility}

\subsection{Limitations of the Base Form}

The base formula has:
\begin{itemize}
    \item a uniform distribution of lobes
    \item no possibility of local clustering
\end{itemize}

To overcome these limitations, we introduce a function $\phi(\theta)$ and define:

\begin{equation}
r(\theta) = \exp\left(-\lambda \left| \sin(k\phi(\theta)) \right|^q \right)
\end{equation}

\subsection{Simple Warp (Without Integrals)}

\begin{equation}
\phi(\theta) = \theta + \mu \sin(\theta - \theta_0)
\end{equation}

Parameters: $\theta_0$ is the cluster position, $\mu$ the deformation
intensity. This simple warp does not guarantee proportional
redistribution and may introduce global distortion.

\subsection{Proportional Reparameterisation (CDF)}

To redistribute angular space without distortion, we define a density:

\begin{equation}
w(t) = 1 + \mu e^{\kappa \cos(t - \theta_0)}
\end{equation}

\begin{equation}
F(\theta) = \int_0^\theta w(t)\,dt
\end{equation}

\begin{equation}
\phi(\theta) = \frac{2\pi}{F(2\pi)} F(\theta)
\end{equation}

\subsection{General REN Formula}

\begin{equation}
r(\theta) = \exp\left(-\lambda \left| \sin(k\phi(\theta)) \right|^q \right)
\end{equation}

\subsection{Interpretation of the Extended Parameters}
\begin{itemize}
    \item $\theta_0$: cluster center
    \item $\mu$: clustering intensity
    \item $\kappa$: cluster width
\end{itemize}

\subsection{Limit Behaviour}
\begin{itemize}
    \item $\mu = 0 \Rightarrow$ classic REN
    \item $\kappa \to \infty \Rightarrow$ point cluster
\end{itemize}

\section{Applications}

\subsection{General Applications}
\begin{itemize}
    \item generation of natural shapes
    \item parametric design
    \item biological modelling
\end{itemize}

\subsection{REN in Hutchinson Synthesis}
\label{sec:hutchinson}

REN was introduced as the geometric constraint used to shape the spatial
dimension of \textbf{Hutchinson Synthesis}, a parametric additive synthesis
system organising sound events over an $N$-dimensional hypervolume of time,
frequency, amplitude and spatial direction \cite{hutchinsonrepo2026}.
This section documents how the base formula above is applied there; a
standalone, dependency-free re-implementation of the same mechanism is
provided in \texttt{../hutchinson\_synthesis\_application/} for reference.

\subsubsection{Role in the Synthesiser}

In Hutchinson Synthesis, azimuth and elevation on the sphere $S^2$ are
treated as two effective dimensions of the sound parameter space, sampled
like any other (frequency, time, amplitude) via the system's stochastic
point-process framework (GIGP). REN is layered on top of that sampling as
an \emph{acceptance-rejection filter}: it does not replace the point
process governing where events tend to occur, it restricts the admissible
region on the sphere to a small number of non-convex ``lobes'', producing
sharp source localisation with near-total exclusion between them.

\subsubsection{From One Curve to One Flower per Niche}

The 1D formula above defines a single boundary curve over the full
angular range $[0, 2\pi)$. In Hutchinson Synthesis it is applied locally,
once per acoustic niche:

\begin{enumerate}
  \item A center $(\text{az}_c, \text{el}_c)$ is placed on the sphere for
  each niche. Centers are spread apart via a repulsive relaxation on
  $S^2$: each point is pushed away from every other point, the resulting
  force is projected onto the local tangent plane (to keep the point on
  the sphere), and the point is renormalised after each step. This
  produces approximately equidistant centers regardless of how many
  niches are active.
  \item Around each center, a local tangent-plane (gnomonic) projection is
  built: any point $(\text{az}, \text{el})$ on the sphere is mapped to
  local polar coordinates $(r, \varphi)$ relative to that center.
  \item The REN boundary is evaluated in this local frame as
  \begin{equation}
    r_{\text{REN}}(\varphi) = R \cdot \exp\!\left(-\lambda
    \left|\sin\!\left(\tfrac{k\varphi}{2}\right)\right|^q\right)
    \label{eq:ren_local}
  \end{equation}
  where $R$ is the niche's radius (auto-scaled from the number of active
  niches, or user-defined). The angle is halved with respect to the base
  formula of eq.~(1) so that the intended number of lobes $k$ is recovered
  after the tangent-plane projection.
\end{enumerate}

\subsubsection{Magnet and Guaranteed Clamp}

Each candidate event position is first assigned to its nearest niche
center (spherical distance), then pulled toward that niche's REN lobes
with a \texttt{force} parameter $\in [0,1]$:

\begin{equation}
r_{\text{final}} = r_{\text{curr}}(1-\text{force}) +
\min(r_{\text{curr}}, r_{\text{REN}}(\varphi))\cdot\text{force}
\end{equation}

At \texttt{force} $=0$ the position is left untouched; at \texttt{force}
$=1$ a final clamp step guarantees the position falls on or inside the
lobes, projecting it onto the nearest lobe border if it would otherwise
fall outside. This is the mechanism behind the REN spatial-distribution
audio and visual examples available in the \texttt{examples/} folder of
the repository \cite{hutchinsonrepo2026}: with high $\lambda$, sources
concentrate tightly in the $k$ lobes, producing near-total spatial silence
in the excluded zones between them.

\subsubsection{A Note on the Codebase}

The Hutchinson Synthesis codebase also contains a separate, exploratory
function (\texttt{ren\_2d}) that combines independent azimuth and
elevation lobe counts multiplicatively. It is not invoked by the
generation pipeline and is not the mechanism described above or used to
produce the published results; it is mentioned here only to avoid
confusion for anyone reading that source directly.

\section{Conclusion}

The REN formula is a minimal yet highly flexible system for generating
complex radial shapes, with direct control over both global and local
structure. Its application inside Hutchinson Synthesis, described above,
shows how the same compact 1D curve can be lifted onto the sphere $S^2$ to
constrain the spatial dimension of a sound synthesiser, without changing
its mathematical core.

\newpage

\begin{thebibliography}{9}

\bibitem{lame}
G. Lam\'e, \textit{Sur les surfaces isothermes}, 1838.

\bibitem{fourier}
J. Fourier, \textit{Th\'eorie analytique de la chaleur}, 1822.

\bibitem{euler}
L. Euler, \textit{Introductio in analysin infinitorum}, 1748.

\bibitem{prob}
W. Feller, \textit{An Introduction to Probability Theory}, Wiley, 1968.

\bibitem{shape}
J. Gielis, ``A generic geometric transformation that unifies a wide range of natural and abstract shapes'', \textit{American Journal of Botany}, 2003.

\bibitem{hutchinsonrepo2026}
A. Di Furia, N. Andreuccetti, \textit{Hutchinson Synthesis} [software repository], \url{https://github.com/anthonydifuria/hutchinson-synthesis}, 2026.

\end{thebibliography}

\end{document}
